% Motor: XeLaTeX. Sesion 16 - material docente UTP.
% Misma estructura docente que la Sesion 15.
\documentclass[aspectratio=169,10pt,c]{beamer}
% Reutiliza sin cambios la identidad visual UTP validada en la Sesion 09.
% Reutiliza sin cambios la identidad visual UTP validada en la Sesion 08.
% Reutiliza sin cambios la identidad visual UTP validada en la Sesion 05.
% La segunda ruta permite compartir los assets comunes sin duplicar binarios.
\input{../../Sesion05/slides/header.tex}
\graphicspath{{../../Sesion05/slides/assets/}}



\graphicspath{{../../Sesion05/slides/assets/}}
% Bloques de código con el mismo estilo que la Sesion 11.
\usepackage{listings}
\lstset{basicstyle=\ttfamily\small,backgroundcolor=\color{utpGrisClaro},
  frame=none,aboveskip=4pt,belowskip=4pt,showstringspaces=false,
  columns=fullflexible,keepspaces=true,breaklines=true,
  xleftmargin=5pt,xrightmargin=5pt,framesep=5pt}


\title{Integración de Servicios Cloud GCP}
\subtitle{Ruta Associate Cloud Engineer: Infraestructura esencial de Google Cloud}
\date{Semana 08 - Sesión 16}
\author{Rel Guzman}

\begin{document}

{\setbeamertemplate{footline}{}\maketitle}

\begin{frame}[c]{Revisión de la sesión anterior}
\begin{block}{La sesión anterior}
\small Conocimos la certificación \textbf{Associate Cloud Engineer} y sus \textbf{4 bloques}, la jerarquía de recursos (\textbf{organización, carpetas, proyectos}), \textbf{regiones y zonas}, la consola, Cloud Shell y \texttt{gcloud}, e hicimos dos labs en Google Skills.
\end{block}
\begin{alertblock}{Pregunta rápida}
\small Creaste una VM en \texttt{us-central1-a}. ¿Esa VM vive en una \textbf{región} o en una \textbf{zona}? ¿Y qué pasa si esa zona se queda sin luz?
\end{alertblock}
\end{frame}

\begin{frame}[c]
\begin{block}{Logro de la sesión}
Al finalizar, sabrás en qué punto de la \textbf{ruta Associate Cloud Engineer} estamos y podrás explicar con tus palabras las \textbf{formas de usar} Google Cloud, cómo funciona su \textbf{red} (subredes, IPs y firewall) y qué decides al crear una \textbf{máquina virtual}.
\end{block}
\end{frame}

\begin{frame}{Saberes previos}
\begin{itemize}\setlength\itemsep{10pt}
  \item En tu casa, el router le da a tu celular una IP como \texttt{192.168.1.25}. ¿Esa IP la puede ver alguien desde Internet?
  \item En la Sesión 14 marcaste \textbf{Permitir tráfico HTTP} al crear la VM. ¿Qué crees que hizo esa casilla por detrás?
  \item Si apagas una VM, ¿crees que Google te sigue cobrando algo?
\end{itemize}
\end{frame}

\begin{frame}[c]{Utilidad de esta sesión}
\begin{center}
\begin{minipage}{.86\textwidth}
{\large Hasta ahora creamos VMs con los valores por defecto. Hoy entiendes \textbf{qué hay detrás}: la red donde vive la VM, quién puede entrar y cuánto cuesta. Es el primer curso técnico de la ruta y la base del bloque 2 del examen.}
\end{minipage}
\end{center}
\end{frame}

% ------------------------------------------------------------- La ruta
\begin{frame}[plain]
\vfill
\begin{center}
{\fontsize{29}{34}\selectfont\bfseries La ruta ACE}\par
\vspace{2mm}
{\small Associate Cloud Engineer: dónde estamos y en qué orden avanzamos}\par
\vspace{4mm}
{\color{utpRojo}\rule{8.5cm}{1.2pt}}
\end{center}
\vfill
\end{frame}

\begin{frame}{La ruta en Google Skills}
\begin{block}{Qué es}
\small La ruta \textbf{Certificación de Associate Cloud Engineer} (\texttt{skills.google/paths/11}) ordena unas \textbf{17 actividades}: cursos, labs e insignias. La seguimos \textbf{en orden}.
\end{block}
\renewcommand{\arraystretch}{1.3}
\centering\small
\begin{tabular}{@{}p{0.6cm}p{6.6cm}p{2.2cm}p{3.2cm}@{}}
\toprule
\textbf{\#} & \textbf{Actividad} & \textbf{Tipo} & \textbf{Hoy}\\
\midrule
1 & Recorrido por los labs prácticos de Google Cloud & Lab & Rápido\\
2 & Crea una guía de estudio para el examen ACE & Curso & \textbf{Lo saltamos}\\
3 & Infraestructura esencial de Google Cloud: conceptos básicos & Curso + labs & \textbf{Lo principal}\\
4+ & Servicios principales, escalamiento, Kubernetes, observabilidad e insignias & Cursos e insignias & Próximas sesiones\\
\bottomrule
\end{tabular}
\par\vspace{2mm}
\begin{alertblock}{Ojo}
\small Google actualiza la ruta seguido: el número y el orden de las actividades pueden cambiar.
\end{alertblock}
\end{frame}

\begin{frame}{Plan de hoy (80 minutos)}
\begin{center}
\begin{tikzpicture}[font=\small,
  paso/.style={draw=utpRojo,rounded corners=4pt,minimum height=1.25cm,align=center}]
  \node[paso,fill=utpRojo!10,minimum width=2.2cm] (a) at (0,0) {\textbf{Recorrido}\\\footnotesize 7 min};
  \node[paso,fill=white,minimum width=2.2cm,dashed] (b) at (2.7,0) {\textbf{Guía ACE}\\\footnotesize 3 min};
  \node[paso,fill=utpRojo!30,minimum width=5.0cm] (c) at (7.0,0) {\textbf{Infraestructura esencial}\\\footnotesize teoría 35 min + práctica 15 min};
  \draw[-latex,thick] (a) -- (b);
  \draw[-latex,thick] (b) -- (c);
\end{tikzpicture}
\end{center}
\begin{columns}[T,onlytextwidth]
\column{.48\textwidth}
\begin{block}{Bloques del curso principal}
\small
\textbf{1.} Interactuar con Google Cloud.\par\smallskip
\textbf{2.} Redes virtuales (VPC).\par\smallskip
\textbf{3.} Máquinas virtuales.
\end{block}
\column{.48\textwidth}
\begin{exampleblock}{El resto del tiempo}
\small Inicio y presentación de la ruta (12 min); preguntas de cierre (8 min).
\end{exampleblock}
\end{columns}
\end{frame}

% ------------------------------------------- Actividad 1: recorrido
\begin{frame}[plain]
\vfill
\begin{center}
{\fontsize{29}{34}\selectfont\bfseries Recorrido por los labs}\par
\vspace{2mm}
{\small Actividad 1: \emph{A Tour of Google Cloud Hands-on Labs} (GSP282), en versión rápida}\par
\vspace{4mm}
{\color{utpRojo}\rule{8.5cm}{1.2pt}}
\end{center}
\vfill
\end{frame}

\begin{frame}{Qué enseña el recorrido}
\begin{block}{Idea}
\small Es el \textbf{primer lab} de la ruta: enseña a usar la plataforma de labs y a moverse por la consola. Ya usamos esa plataforma en las Sesiones 14 y 15.
\end{block}
\begin{columns}[T,onlytextwidth]
\column{.48\textwidth}
\begin{exampleblock}{La plataforma de labs}
\small
\textbf{Start Lab} crea un \textbf{proyecto temporal}.\par\smallskip
Te da un \textbf{usuario y contraseña} de prueba y un \textbf{temporizador}.\par\smallskip
\textbf{Check my progress} califica cada tarea.\par\smallskip
\textbf{End Lab} borra todo al terminar.
\end{exampleblock}
\column{.48\textwidth}
\begin{block}{Lo que muestra de Google Cloud}
\small
\textbf{Proyectos:} nombre, ID y número.\par\smallskip
\textbf{IAM:} usuarios, roles y permisos.\par\smallskip
\textbf{APIs:} buscar y habilitar una API en la biblioteca.
\end{block}
\end{columns}
\par\vspace{1mm}
\begin{alertblock}{La regla de oro}
\small Abre la consola en \textbf{incógnito} con el usuario temporal, nunca con tu cuenta personal.
\end{alertblock}
\end{frame}

\begin{frame}{¿Por qué lo vemos rápido?}
\begin{columns}[T,onlytextwidth]
\column{.48\textwidth}
\begin{alertblock}{Lo desactualizado}
\small
Las \textbf{capturas} son de una consola anterior: los menús ya no se ven igual.\par\smallskip
Habla de \textbf{Qwiklabs} y \textbf{Cloud Skills Boost}: hoy la plataforma se llama \textbf{Google Skills}.\par\smallskip
Algunas rutas de menú cambiaron de nombre o de lugar.
\end{alertblock}
\column{.48\textwidth}
\begin{exampleblock}{Lo que sigue valiendo}
\small
Proyecto, IAM y APIs: \textbf{los conceptos no cambian}.\par\smallskip
Si un menú no aparece, usa la \textbf{barra de búsqueda} de la consola: escribe \emph{IAM} o \emph{API Library}.
\end{exampleblock}
\begin{block}{Qué hacemos}
\small Lo recorro en pantalla en pocos minutos. Si quieres que figure como completado en la ruta, hazlo de tarea.
\end{block}
\end{columns}
\end{frame}

% --------------------------------------------- Actividad 2: guia ACE
\begin{frame}[plain]
\vfill
\begin{center}
{\fontsize{29}{34}\selectfont\bfseries Guía de estudio ACE}\par
\vspace{2mm}
{\small Actividad 2: \emph{Preparing for Your Associate Cloud Engineer Journey}}\par
\vspace{4mm}
{\color{utpRojo}\rule{8.5cm}{1.2pt}}
\end{center}
\vfill
\end{frame}

\begin{frame}{La guía de estudio: la saltamos, y por qué}
\begin{block}{Qué es}
\small Un curso para \textbf{planificar} tu preparación: recorre los 4 bloques del examen con \textbf{preguntas de diagnóstico} y al final armas tu propio plan de estudio.
\end{block}
\begin{columns}[T,onlytextwidth]
\column{.48\textwidth}
\begin{alertblock}{Por qué no ahora}
\small
\textbf{1.} Las preguntas de diagnóstico suponen que \textbf{ya estudiaste}: hoy fallarías casi todas sin aprender nada nuevo.\par\smallskip
\textbf{2.} No tiene \textbf{labs}: no practicas en la consola.\par\smallskip
\textbf{3.} Lo esencial (bloques del examen y plan de estudio) ya lo vimos en la \textbf{Sesión 15}.
\end{alertblock}
\column{.48\textwidth}
\begin{exampleblock}{Cuándo volver}
\small Al \textbf{terminar la ruta}, antes de inscribirte al examen. Ahí sí sirve para encontrar tus puntos débiles.
\end{exampleblock}
\begin{block}{Una comparación}
\small Es como un \textbf{simulacro} de admisión: se da cuando ya estudiaste, no el primer día de academia.
\end{block}
\end{columns}
\end{frame}

% --------------------------------- Actividad 3: infraestructura esencial
\begin{frame}[plain]
\vfill
\begin{center}
{\fontsize{29}{34}\selectfont\bfseries Infraestructura esencial}\par
\vspace{2mm}
{\small Actividad 3: \emph{Essential Google Cloud Infrastructure: Foundation} (conceptos básicos)}\par
\vspace{4mm}
{\color{utpRojo}\rule{8.5cm}{1.2pt}}
\end{center}
\vfill
\end{frame}

\begin{frame}{Mapa del curso}
\begin{block}{Qué es}
\small El \textbf{primer curso técnico} de la ruta (unas 7 horas). Se centra en \textbf{Compute Engine}: la red donde viven las VMs y las VMs mismas.
\end{block}
\renewcommand{\arraystretch}{1.3}
\centering\small
\begin{tabular}{@{}p{4.8cm}p{7.8cm}@{}}
\toprule
\textbf{Módulo} & \textbf{Labs, en orden}\\
\midrule
1. Interactuar con Google Cloud & Consola y Cloud Shell; una web lista con Marketplace\\
2. Redes virtuales & Crear redes VPC; VMs sin IP externa con Cloud NAT\\
3. Máquinas virtuales & Crear VMs; configurar VMs\\
\bottomrule
\end{tabular}
\par\vspace{3mm}
\begin{alertblock}{Cómo lo trabajamos}
\small En clase vemos la \textbf{teoría} de los tres módulos. Los labs van en orden: hoy empiezas el primero.
\end{alertblock}
\end{frame}

% ---- Modulo 1
\begin{frame}[fragile]{Módulo 1: cuatro formas de hablar con Google Cloud}
\begin{block}{Idea}
\small Google Cloud es como un banco: puedes ir a la \textbf{ventanilla}, usar el \textbf{cajero}, la \textbf{app} o que tu sistema se conecte \textbf{directo}. El dinero es el mismo; cambia la forma.
\end{block}
\begin{columns}[T,onlytextwidth]
\column{.235\textwidth}
\begin{exampleblock}{Consola}
\footnotesize Página web con clics. Ideal para \textbf{ver} y empezar.
\end{exampleblock}
\column{.235\textwidth}
\begin{exampleblock}{Cloud Shell}
\footnotesize Una terminal en el navegador, ya lista. No instalas nada.
\end{exampleblock}
\column{.235\textwidth}
\begin{exampleblock}{\texttt{gcloud}}
\footnotesize Comandos que puedes instalar en tu laptop o en WSL.
\end{exampleblock}
\column{.235\textwidth}
\begin{exampleblock}{API}
\footnotesize Tus programas (Python, Node.js) le piden cosas a Google Cloud.
\end{exampleblock}
\end{columns}
\par\vspace{2mm}
\begin{block}{Lo mismo, de dos formas}
\small Con clics: \textbf{Compute Engine} $\rightarrow$ \textbf{Instancias de VM}. Con un comando:
\begin{lstlisting}[basicstyle=\ttfamily\footnotesize]
gcloud compute instances list
\end{lstlisting}
\end{block}
\end{frame}

\begin{frame}{Módulo 1: el proyecto y el Marketplace}
\begin{columns}[T,onlytextwidth]
\column{.48\textwidth}
\begin{block}{Un proyecto tiene tres datos}
\small
\textbf{Nombre:} como tu apodo; lo cambias cuando quieras.\par\smallskip
\textbf{ID:} como tu \textbf{DNI}; único y \textbf{no cambia nunca}. Es el que piden los comandos.\par\smallskip
\textbf{Número:} lo pone Google.
\end{block}
\begin{alertblock}{Error típico}
\small Escribir el \textbf{nombre} cuando el comando pide el \textbf{ID}.
\end{alertblock}
\column{.48\textwidth}
\begin{exampleblock}{Marketplace: soluciones listas}
\small Una \textbf{tienda} de programas ya armados: WordPress, una tienda online, un servidor web con base de datos. Eliges uno y Google crea la VM y todo lo necesario.
\end{exampleblock}
\begin{block}{Una comparación}
\small Es comprar un \textbf{mueble armado} en vez de las piezas: rápido, pero igual debes saber qué compraste y cuánto cuesta.
\end{block}
\end{columns}
\end{frame}

% ---- Modulo 2
\begin{frame}{Módulo 2: ¿qué es una red VPC?}
\begin{block}{Definición}
\small La \textbf{red privada de tu proyecto} (\emph{Virtual Private Cloud}). Es donde se conectan tus VMs, como el wifi de tu casa conecta tus equipos.
\end{block}
\begin{columns}[c,onlytextwidth]
\column{.52\textwidth}
\centering
\begin{tikzpicture}[font=\footnotesize]
  \draw[utpRojo,thick,rounded corners=6pt,fill=utpRojo!5] (0,0) rectangle (6.4,3.3);
  \node[anchor=north west,font=\small\bfseries] at (0.1,3.25) {Red VPC: el condominio};
  \draw[rounded corners=4pt,fill=white] (0.3,0.3) rectangle (3.0,2.6);
  \node[anchor=north west] at (0.35,2.55) {\textbf{Subred en EE.\,UU.}};
  \node[anchor=north west] at (0.35,2.15) {un bloque};
  \node[draw,fill=utpRojo!20,rounded corners=2pt] at (1.0,0.9) {VM a};
  \node[draw,fill=utpRojo!20,rounded corners=2pt] at (2.3,0.9) {VM b};
  \draw[rounded corners=4pt,fill=white] (3.4,0.3) rectangle (6.1,2.6);
  \node[anchor=north west] at (3.45,2.55) {\textbf{Subred en Chile}};
  \node[anchor=north west] at (3.45,2.15) {otro bloque};
  \node[draw,fill=utpRojo!20,rounded corners=2pt] at (4.75,0.9) {VM c};
\end{tikzpicture}
\column{.44\textwidth}
\begin{exampleblock}{Una comparación}
\small La VPC es un \textbf{condominio}; las \textbf{subredes} son los bloques de edificios; las \textbf{VMs}, los departamentos.
\end{exampleblock}
\begin{block}{Lo sorprendente}
\small Cada subred está en \textbf{una región}, pero la red las une a todas: la VM de Chile y la de EE.\,UU. se hablan como vecinas, sin salir a Internet.
\end{block}
\end{columns}
\par\vspace{1mm}
\begin{alertblock}{Ya la usaste}
\small Cada proyecto trae una red llamada \texttt{default}. Tus VMs de la Sesión 14 vivían ahí.
\end{alertblock}
\end{frame}

\begin{frame}{Módulo 2: subredes y sus direcciones}
\begin{block}{Cada subred tiene un rango de direcciones}
\small Se escribe así: \texttt{10.0.1.0/24}. El número después de la barra dice el \textbf{tamaño}: mientras \textbf{más chico}, \textbf{más direcciones}.
\end{block}
\renewcommand{\arraystretch}{1.3}
\centering\small
\begin{tabular}{@{}p{2.4cm}p{3.4cm}p{3.0cm}p{3.8cm}@{}}
\toprule
\textbf{Rango} & \textbf{Direcciones} & \textbf{Para VMs} & \textbf{Alcanza para\ldots}\\
\midrule
\texttt{/29} & 8 & 4 & una prueba pequeña\\
\texttt{/24} & 256 & 252 & un proyecto normal\\
\texttt{/20} & 4\,096 & 4\,092 & una empresa grande\\
\bottomrule
\end{tabular}
\par\vspace{3mm}
\begin{columns}[T,onlytextwidth]
\column{.48\textwidth}
\begin{exampleblock}{¿Por qué 4 menos?}
\small Google se \textbf{reserva 4 direcciones} en cada subred para que la red funcione (la primera, la puerta de salida y las dos últimas).
\end{exampleblock}
\column{.48\textwidth}
\begin{alertblock}{Si te quedas corto}
\small Puedes \textbf{agrandar} la subred sin apagar nada. Lo que no se puede es achicarla.
\end{alertblock}
\end{columns}
\end{frame}

\begin{frame}{Módulo 2: IP interna e IP externa}
\begin{columns}[T,onlytextwidth]
\column{.48\textwidth}
\begin{block}{IP interna}
\small La tiene \textbf{toda} VM. Sirve para hablar \textbf{dentro} de la red, con las otras VMs.
\end{block}
\begin{block}{IP externa}
\small \textbf{Opcional}. Es la que abriste desde el navegador en el lab de NGINX.\par\smallskip
\textbf{Efímera:} puede cambiar si apagas la VM.\par
\textbf{Estática:} la reservas y se queda fija.
\end{block}
\column{.48\textwidth}
\begin{exampleblock}{Una comparación}
\small La IP interna es el \textbf{anexo} de una oficina: solo funciona dentro. La externa es el \textbf{número público} de la central: cualquiera puede llamar.
\end{exampleblock}
\begin{alertblock}{Saberes previos, resuelto}
\small El \texttt{192.168.1.25} de tu celular es como una IP \textbf{interna}: desde Internet no se ve. Tu casa sale con la IP pública del router.
\end{alertblock}
\end{columns}
\end{frame}

\begin{frame}{Módulo 2: el firewall, el portero de la red}
\begin{center}
\begin{tikzpicture}[font=\small,
  caja/.style={draw=utpRojo,rounded corners=4pt,minimum width=2.4cm,minimum height=1.1cm,align=center}]
  \node[caja,fill=white] (i) at (0,0) {\textbf{Internet}};
  \node[caja,fill=utpRojo!30] (f) at (5.2,0) {\textbf{Firewall}\\\footnotesize revisa su lista};
  \node[caja,fill=utpRojo!10] (v) at (10.4,0) {\textbf{Tu VM}};
  \draw[-latex,thick] (i) -- node[above,font=\footnotesize] {visita al puerto 80} (f);
  \draw[-latex,thick] (f) -- node[above,font=\footnotesize] {¿está en la lista?} (v);
\end{tikzpicture}
\end{center}
\begin{columns}[T,onlytextwidth]
\column{.48\textwidth}
\begin{block}{La regla básica}
\small Lo que \textbf{entra}: todo \textbf{bloqueado}, salvo lo que tú permitas.\par\smallskip
Lo que \textbf{sale}: todo permitido.
\end{block}
\begin{exampleblock}{Una regla dice}
\small \textbf{Qué} (puerto 80, web), \textbf{desde dónde} (cualquier IP) y \textbf{a quién} (las VMs con cierta etiqueta).
\end{exampleblock}
\column{.48\textwidth}
\begin{block}{Los puertos más comunes}
\small
\textbf{22:} SSH, entrar a la terminal.\par
\textbf{80:} web (\texttt{http://}).\par
\textbf{443:} web segura (\texttt{https://}).
\end{block}
\begin{alertblock}{Una comparación}
\small El portero del condominio con su \textbf{lista de invitados}: si no estás, no entras.
\end{alertblock}
\end{columns}
\end{frame}

\begin{frame}[fragile]{Módulo 2: la casilla ``Permitir tráfico HTTP''}
\begin{block}{Qué hizo en la Sesión 14}
\small Le puso a la VM una \textbf{etiqueta}: \texttt{http-server}. La red ya tenía una regla: ``a las VMs con esa etiqueta, deja pasar la web (puerto 80)''.
\end{block}
\begin{center}
\begin{tikzpicture}[font=\small,
  caja/.style={draw=utpRojo,rounded corners=4pt,minimum width=3.6cm,minimum height=1.1cm,align=center}]
  \node[caja,fill=utpRojo!10] (c) at (0,0) {Marcas la casilla};
  \node[caja,fill=utpRojo!20] (e) at (4.6,0) {VM con etiqueta\\\texttt{http-server}};
  \node[caja,fill=utpRojo!35] (r) at (9.2,0) {La regla la deja\\pasar: la web abre};
  \draw[-latex,thick] (c) -- (e);
  \draw[-latex,thick] (e) -- (r);
\end{tikzpicture}
\end{center}
\begin{exampleblock}{La misma regla, escrita como comando}
\begin{lstlisting}[basicstyle=\ttfamily\footnotesize]
gcloud compute firewall-rules create permitir-web \
  --allow=tcp:80 --target-tags=http-server
\end{lstlisting}
\end{exampleblock}
\begin{alertblock}{Error típico}
\small Crear la regla pero olvidar ponerle la \textbf{etiqueta} a la VM: la página no abre.
\end{alertblock}
\end{frame}

\begin{frame}{Módulo 2: VMs sin IP externa y lo que cuesta la red}
\begin{columns}[T,onlytextwidth]
\column{.48\textwidth}
\begin{block}{VMs escondidas, pero útiles}
\small Una VM \textbf{sin IP externa} es más segura: nadie la alcanza desde Internet. Pero, ¿cómo descarga actualizaciones?\par\smallskip
\textbf{Cloud NAT} le da \textbf{salida} a Internet, sin dejar que nadie entre.
\end{block}
\begin{exampleblock}{Una comparación}
\small Como una casa con \textbf{puerta solo de salida}: tú sales a comprar, pero nadie toca el timbre.
\end{exampleblock}
\column{.48\textwidth}
\begin{alertblock}{Qué se cobra}
\small
Datos que \textbf{entran}: gratis.\par\smallskip
Datos entre VMs \textbf{cercanas} (misma zona): gratis.\par\smallskip
Datos que \textbf{salen a Internet}: se cobran.
\end{alertblock}
\begin{block}{Idea clave}
\small Lo que \textbf{sale} cuesta. Mantener las cosas cerca ahorra dinero.
\end{block}
\end{columns}
\end{frame}

% ---- Modulo 3
\begin{frame}{Módulo 3: una VM por dentro}
\begin{block}{Qué eliges al crear una VM}
\small Lo mismo que al comprar una laptop: \textbf{procesador} (vCPU), \textbf{memoria}, \textbf{disco} y \textbf{sistema operativo}. Más \textbf{dónde vive} (zona) y \textbf{a qué red} se conecta.
\end{block}
\renewcommand{\arraystretch}{1.3}
\centering\small
\begin{tabular}{@{}p{3.6cm}p{2.8cm}p{6.2cm}@{}}
\toprule
\textbf{Familia} & \textbf{Ejemplo} & \textbf{Como una\ldots}\\
\midrule
Uso general & E2, N2 & \textbf{laptop de oficina}: sirve para casi todo\\
Mucho cálculo & C2 & \textbf{PC gamer}: procesador potente\\
Mucha memoria & M3 & \textbf{servidor gigante} para bases de datos enormes\\
Con GPU & A3, G2 & \textbf{equipo de IA}: tarjetas gráficas\\
\bottomrule
\end{tabular}
\par\vspace{3mm}
\begin{columns}[T,onlytextwidth]
\column{.48\textwidth}
\begin{exampleblock}{Cómo se lee \texttt{e2-medium}}
\small Familia \textbf{E2}, tamaño \textbf{medium}. Es la que usamos en los labs.
\end{exampleblock}
\column{.48\textwidth}
\begin{alertblock}{Para empezar}
\small Casi siempre: \textbf{E2}. Es la más barata.
\end{alertblock}
\end{columns}
\end{frame}

\begin{frame}{Módulo 3: encendida, detenida, borrada}
\begin{center}
\begin{tikzpicture}[font=\small,
  est/.style={draw=utpRojo,rounded corners=4pt,minimum width=3.4cm,minimum height=1.3cm,align=center}]
  \node[est,fill=utpRojo!40] (r) at (0,0) {\textbf{Encendida}\\\footnotesize\texttt{RUNNING}};
  \node[est,fill=utpRojo!12] (t) at (4.8,0) {\textbf{Detenida}\\\footnotesize\texttt{TERMINATED}};
  \node[est,fill=white] (b) at (9.6,0) {\textbf{Borrada}\\\footnotesize ya no existe};
  \draw[-latex,thick] ([yshift=2mm]r.east) -- node[above,font=\footnotesize] {detener} ([yshift=2mm]t.west);
  \draw[-latex,thick] ([yshift=-2mm]t.west) -- node[below,font=\footnotesize] {iniciar} ([yshift=-2mm]r.east);
  \draw[-latex,thick] (t) -- node[above,font=\footnotesize] {borrar} (b);
\end{tikzpicture}
\end{center}
\renewcommand{\arraystretch}{1.3}
\centering\small
\begin{tabular}{@{}p{3.6cm}p{2.6cm}p{2.6cm}p{2.6cm}@{}}
\toprule
\textbf{¿Se cobra?} & \textbf{Encendida} & \textbf{Detenida} & \textbf{Borrada}\\
\midrule
Procesador y memoria & Sí & No & No\\
Disco & Sí & \textbf{Sí} & No\\
\bottomrule
\end{tabular}
\par\vspace{3mm}
\begin{alertblock}{Ojo con el nombre}
\small \emph{Terminated} suena a ``borrada'', pero solo está \textbf{apagada}. El disco se sigue cobrando: al terminar un proyecto, \textbf{bórrala}.
\end{alertblock}
\end{frame}

\begin{frame}{Módulo 3: cómo pagar menos por una VM}
\renewcommand{\arraystretch}{1.35}
\centering\small
\begin{tabular}{@{}p{3.4cm}p{5.2cm}p{4.0cm}@{}}
\toprule
\textbf{Opción} & \textbf{Cómo funciona} & \textbf{Como un\ldots}\\
\midrule
A pedido & Pagas por segundo mientras está encendida & \textbf{taxi} con taxímetro\\
Uso comprometido & Prometes usarla 1 o 3 años y te hacen descuento & \textbf{plan anual} del gimnasio\\
\textbf{VM Spot} & Muy barata, pero Google puede apagarla cuando necesite espacio & asiento en \textbf{lista de espera}\\
\bottomrule
\end{tabular}
\par\vspace{4mm}
\begin{columns}[T,onlytextwidth]
\column{.48\textwidth}
\begin{exampleblock}{Cuándo usar Spot}
\small Trabajos que se pueden \textbf{repetir} si se cortan: procesar fotos por la noche, cálculos largos.
\end{exampleblock}
\column{.48\textwidth}
\begin{alertblock}{Cuándo no}
\small La web de tu tienda: si Google la apaga, tus clientes no compran.
\end{alertblock}
\end{columns}
\end{frame}

\begin{frame}{Módulo 3: discos, imágenes y copias}
\begin{columns}[T,onlytextwidth]
\column{.31\textwidth}
\begin{block}{Imagen}
\small El \textbf{sistema operativo} con el que nace la VM: Debian, Ubuntu, Windows\ldots{}\par\smallskip
También puedes crear tu \textbf{propia imagen} de una VM ya configurada.
\end{block}
\column{.31\textwidth}
\begin{block}{Disco}
\small Donde se guardan los archivos y el sistema operativo.\par\smallskip
Se puede \textbf{agrandar}, no achicar.
\end{block}
\column{.31\textwidth}
\begin{block}{Instantánea}
\small Una \textbf{foto} del disco en un momento. Sirve como copia de seguridad o para clonar la VM.
\end{block}
\end{columns}
\par\vspace{3mm}
\begin{exampleblock}{Una comparación}
\small La \textbf{imagen} es el molde de la torta; el \textbf{disco}, la torta que te comes; la \textbf{instantánea}, la foto que le tomas antes de cortarla.
\end{exampleblock}
\begin{alertblock}{Script de inicio}
\small Un pequeño programa que la VM corre \textbf{cada vez que arranca}. Así se puede instalar NGINX sin entrar por SSH.
\end{alertblock}
\end{frame}

% ------------------------------------------------------------ Practica
\begin{frame}[plain]
\vfill
\begin{center}
{\fontsize{29}{34}\selectfont\bfseries Práctica}\par
\vspace{2mm}
{\small Los labs del curso}\par
\vspace{4mm}
{\color{utpRojo}\rule{8.5cm}{1.2pt}}
\end{center}
\vfill
\end{frame}

\begin{frame}{Práctica: redes y VMs \textnormal{\small(15 min en clase + tarea)}}
\begin{block}{En clase}
\small
1. Calcula en tu cuaderno cuántas IPs quedan para VMs en una subred \texttt{/24}, una \texttt{/20} y una \texttt{/29}. Recuerda las 4 que reserva Google.\par\smallskip
2. Escribe qué reglas de firewall dejarían pasar el tráfico web pero no el SSH desde Internet.\par\smallskip
3. Inicia el primer lab del curso: \emph{Explore the Google Cloud Console and Cloud Shell}.
\end{block}
\begin{columns}[T,onlytextwidth]
\column{.48\textwidth}
\begin{exampleblock}{De tarea, en orden}
\small Terminar el de consola y Cloud Shell, luego \emph{Create and Configure VPC Networks} y \emph{Create Virtual Machines}. Los otros tres del curso, opcionales.
\end{exampleblock}
\column{.48\textwidth}
\begin{alertblock}{Entregable}
\small Captura de cada lab con todos los \textbf{Check my progress} en verde y tus cálculos de subredes.
\end{alertblock}
\end{columns}
\end{frame}

\begin{frame}{Cierre: preguntas para comprobar}
\begin{block}{Responde con tus palabras}
\small
1. ¿Por qué saltamos la guía de estudio del examen ACE y cuándo conviene hacerla?\par\medskip
2. Una red VPC, ¿es global, regional o zonal? ¿Y una subred?\par\medskip
3. ¿Cuántas IPs quedan para VMs en una subred \texttt{/24}? ¿Por qué no son 256?\par\medskip
4. Creaste una regla que abre el puerto 80 para la etiqueta \texttt{web}, pero la página no abre. ¿Qué revisas?\par\medskip
5. Apagaste tu VM. ¿Qué te siguen cobrando?\par\medskip
6. Debes procesar miles de fotos por la noche al menor costo. ¿Qué tipo de VM usas?
\end{block}
\begin{alertblock}{Lo que viene}
\small Seguimos la ruta: \textbf{Infraestructura esencial: servicios principales} (IAM a fondo, almacenamiento y bases de datos).
\end{alertblock}
\end{frame}

\begin{frame}[plain]
\begin{tikzpicture}[remember picture,overlay]
  \node[anchor=center,inner sep=0pt] at (current page.center)
    {\includegraphics[width=\paperwidth,height=\paperheight]{gracias-final.pdf}};
\end{tikzpicture}
\end{frame}

\begin{frame}{Referencias}
\scriptsize
\begin{itemize}\setlength\itemsep{3pt}
  \item Google Skills: ruta \emph{Certificación de Associate Cloud Engineer}, \texttt{skills.google/paths/11}.
  \item Google Skills: \emph{A Tour of Google Cloud Hands-on Labs} (GSP282) y \emph{Preparing for Your Associate Cloud Engineer Journey}.
  \item Google Skills: \emph{Essential Google Cloud Infrastructure: Foundation}, \texttt{skills.google/paths/11/course\_templates/50}.
  \item Google Cloud: \emph{VPC network overview}, \emph{Subnets} y \emph{VPC firewall rules}, \texttt{cloud.google.com/vpc/docs}.
  \item Google Cloud: \emph{Cloud NAT overview} y \emph{Private Google Access}, \texttt{cloud.google.com/nat/docs} y \texttt{cloud.google.com/vpc/docs}.
  \item Google Cloud: \emph{Machine families resource and comparison guide}, \emph{VM instance lifecycle} y \emph{Spot VMs}, \texttt{cloud.google.com/compute/docs}.
  \item Google Cloud: \emph{Disks and images}, \texttt{cloud.google.com/compute/docs/disks}.
  \item Google Cloud: \emph{Network pricing}, \texttt{cloud.google.com/vpc/network-pricing}.
  \item Material de repaso: \texttt{Sesion15/slides/slides.tex}.
  \item Las respuestas de todas las preguntas de esta sesión están en \texttt{respuestas.md}.
\end{itemize}
\end{frame}

\end{document}
